% BHU PYQ -- Manually transcribed & error-corrected
% Paper: MTB-202 : Statics and Dynamics
% Exam: B.A. (Hons.) Semester II Examination, 2023-24

\documentclass[11pt,a4paper]{article}
\usepackage[a4paper,top=2.5cm,bottom=2.5cm,left=2.8cm,right=2.8cm,headheight=14pt]{geometry}
\usepackage{amsmath,amssymb}
\usepackage[T1]{fontenc}
\usepackage[utf8]{inputenc}
\usepackage{lmodern}
\usepackage{microtype}
\usepackage{fancyhdr}
\usepackage{enumitem}
\usepackage{hyperref}
\usepackage{lastpage}
\usepackage{parskip}

\hypersetup{colorlinks=true,urlcolor=black,linkcolor=black,
  pdftitle={MTB-202: Statics and Dynamics -- BHU PYQ},pdfauthor={Banaras Hindu University}}

\pagestyle{fancy}\fancyhf{}
\renewcommand{\headrulewidth}{0.4pt}\renewcommand{\footrulewidth}{0.4pt}
\fancyhead[L]{\small\textit{Banaras Hindu University}}
\fancyhead[R]{\small\textit{MTB-202: Statics and Dynamics}}
\fancyfoot[C]{\small Page \thepage\ of \pageref{LastPage}}

\newlist{parts}{enumerate}{2}
\setlist[parts,1]{label=(\alph*),leftmargin=*,itemsep=5pt,topsep=3pt}
\setlist[parts,2]{label=(\roman*),leftmargin=*,itemsep=3pt,topsep=2pt}
\newcommand{\pts}[1]{\hfill\textit{[#1]}}

\begin{document}

\begin{center}
  {\large\bfseries BANARAS HINDU UNIVERSITY}\\[2pt]
  {\normalsize B.A. (Hons.) --- Semester II Examination, 2023-24}\\[6pt]
  \rule{0.8\linewidth}{0.5pt}\\[6pt]
  {\large\bfseries Paper No. MTB-202: Statics and Dynamics}\\[4pt]
  {\normalsize Time: 3 Hours \hfill Full Marks: 70}
\end{center}
\rule{\linewidth}{0.4pt}
\smallskip
\noindent\textit{Instructions: Answer \textbf{five} questions including Question~1 which is compulsory. Symbols have their usual meanings.}
\bigskip

\noindent\textbf{Question 1.} \textit{(Compulsory --- Answer all parts. 2 marks each.)}
\begin{parts}
  \item Determine whether the forces $\vec{F}_1 = -2\vec{i} + 3\vec{j}$, $\vec{F}_2 = 4\vec{i} - 4\vec{j}$, $\vec{F}_3 = -3\vec{i} - \vec{j}$ and $\vec{F}_4 = \vec{i} + 2\vec{j}$ acting at the points $(2, 1)$, $(-1, 2)$, $(1, -1)$ and $(-1, 0)$ respectively are in equilibrium or not. \pts{2}
  \item Find the virtual work done by the thrust of a rod. \pts{2}
  \item State Kepler's laws of planetary motion. \pts{2}
  \item A particle describes the parabola $p^2 = ar$ under a force which is always directed towards its focus. Find the law of force. \pts{2}
  \item What is a cycloidal pendulum? \pts{2}
  \item Show that the areal velocity is constant for a particle moving in a central orbit. \pts{2}
  \item Prove that the relation between angular velocity $\dot{\theta}$ and linear velocity $v$ is:
  \[
  \dot{\theta} = \frac{vp}{r^2}
  \]
  where $\dot{\theta}$ is the angular velocity, $v$ is the linear velocity, $r$ is the radius vector, and $p$ is the perpendicular distance from the origin on the tangent. \pts{2}
\end{parts}

\medskip\rule{\linewidth}{0.2pt}

\medskip\noindent\textbf{Question 2.}
\begin{parts}
  \item Show that a system of coplanar forces acting at different points of a rigid body can be reduced to a single force through a given point and a couple. \pts{7}
  \item Forces of relative magnitudes 5, 1, 1, and $-3$ act along the sides $AB$, $BC$, $CD$, and $DA$ respectively of a square $ABCD$. Show that the system is equivalent to a single force and find its magnitude, direction, and line of action. \pts{7}
\end{parts}

\medskip\rule{\linewidth}{0.2pt}

\medskip\noindent\textbf{Question 3.}
\begin{parts}
  \item State and prove the principle of virtual work for a system of coplanar forces acting on a particle. \pts{7}
  \item A heavy rod of length $2l$ rests upon a fixed smooth peg at $C$ and with its end $B$ upon a smooth curve. If it rests in all positions, show that the curve is a conchoid whose polar equation with $C$ as origin (pole) is:
  \[
  r = a + \frac{l}{\sin\theta}
  \] \pts{7}
\end{parts}

\medskip\rule{\linewidth}{0.2pt}

\medskip\noindent\textbf{Question 4.}
\begin{parts}
  \item Derive the radial and transverse components of acceleration of a particle moving in a plane curve. \pts{7}
  \item If the radial and transverse velocities of a point are always proportional to each other and this holds for the accelerations also, prove that its velocity will vary as some power of the radius vector. \pts{7}
\end{parts}

\medskip\rule{\linewidth}{0.2pt}

\medskip\noindent\textbf{Question 5.}
\begin{parts}
  \item If the central orbit is $r^n = a^n \cos n\theta$ under a force towards the pole, then find the law of force and the velocity at any point. \pts{7}
  \item A particle moves in a plane with an acceleration $F$ which is directed towards a fixed point $O$ in the plane. Find the differential equation of its path. \pts{7}
\end{parts}

\medskip\rule{\linewidth}{0.2pt}

\medskip\noindent\textbf{Question 6.}
\begin{parts}
  \item If in a simple harmonic motion $u, v, w$ be the velocities at distances $a, b, c$ from a fixed point on the straight line which is not the centre of the force, show that the period $T$ is given by the equation:
  \[
  \frac{4\pi^2}{T^2}(b-c)(c-a)(a-b) = \begin{vmatrix} u^2 & v^2 & w^2 \\ a & b & c \\ 1 & 1 & 1 \end{vmatrix}
  \] \pts{7}
  \item A particle attached to a fixed peg $O$ by a string of length $l$ is lifted up with the string horizontal and then let go. Prove that when the string makes an angle $\theta$ with the horizontal, the resultant acceleration is:
  \[
  g \sqrt{1 + 3\sin^2\theta}
  \] \pts{7}
\end{parts}

\medskip\rule{\linewidth}{0.2pt}

\medskip\noindent\textbf{Question 7.}
\begin{parts}
  \item A particle slides down a smooth cycloid whose axis is vertical and vertex downwards. Find the velocity of the particle and the reaction on it at any point of the cycloid. \pts{7}
  \item Obtain the equations of motion of a particle in a frame of reference moving with translational motion and rotating about the common origin with respect to a fixed frame. \pts{7}
\end{parts}

\vfill
\rule{\linewidth}{0.4pt}
\begin{center}\small
  \href{https://bhu-exam.vercel.app/}{Click here to visit our website}\\[4pt]
  \href{https://me-aryan.vercel.app/}{Click here to visit our contributor}\\[4pt]
  \href{https://quashan-me.vercel.app/}{Click here to visit our contributor}
\end{center}

\end{document}
